\section{Outline of the Dissertation}

This dissertation is organized into seven chapters.
Chapter~1 is the introduction, which motivates the thesis theme of
\emph{defending emerging software systems in untrusted environments},
reviews the most relevant prior work,
and explains how the four technical works of the dissertation fit together
under one defense perspective.

\noindent \textbf{Chapter 2: Preliminaries and Unified Notation.}
This chapter introduces the technical background and preliminaries 
used throughout the dissertation.
The chapter then fixes the notation reused across the later technical chapters.

\noindent \textbf{Chapter 3: Defending Smart Contracts with Runtime Invariants.}
This chapter studies runtime invariant generation and enforcement for smart-contract security.
It is based on the observation that malicious transactions often exhibit abnormal behaviors
in asset movements, state updates, or invocation patterns when compared with benign usage.
Building on this observation, the chapter presents techniques for learning behavioral invariants
from historical transaction traces and enforcing them as runtime guards for deployed contracts.
The material in this chapter is based on the published paper
\textsc{Trace2Inv}~\cite{chen2024demystifying}.

\noindent \textbf{Chapter 4: Defending Smart Contracts with Control-Flow Integrity.}
This chapter studies control-flow integrity for decentralized finance smart contracts.
Its central observation is that many attacks exploit unexpected control flows
across functions and contracts, even when each individual component may appear benign in isolation.
Accordingly, the chapter develops a runtime guard framework that instruments contracts. pre-deployment
and validates whether transaction executions follow legitimate control-flow structures.
The material in this chapter is based on the accepted paper
\textsc{CrossGuard}~\cite{chen2025secure}.

\noindent \textbf{Chapter 5: Auditing Malicious Artifacts in Production LLMs.}
This chapter shifts from on-chain runtime protection to upstream risk in AI-assisted software development.
It studies whether production large language models can generate code containing malicious artifacts,
such as scam or phishing endpoints, under ordinary developer-style prompts.
To answer this question, the chapter presents an automated auditing framework
that systematically constructs benign-looking coding requests and measures malicious code generation rates
in production models.
The material in this chapter is based on the written paper
\textsc{Scam2Prompt}~\cite{chen2025scam2prompt}, which is currently under submission.

\noindent \textbf{Chapter 6: Defending Coding Agents via Neutralizing External Poisons.}
This chapter presents the planned final work of the dissertation.
It studies a more upstream threat model in which coding agents rely on external information
during task solving.
The chapter examines how external content can be retrieved, trusted,
and translated into harmful coding outcomes, and it proposes defenses that reduce this risk
while preserving normal agent utility. 
The material in this chapter is based on ongoing research.

\noindent \textbf{Chapter 7: Conclusion.}
This chapter concludes the dissertation by summarizing the four defense layers,
revisiting the common theme of defending emerging software systems in untrusted environments,
and discussing broader lessons for securing increasingly autonomous software systems.
